\section{Introduction}

An agent helping a user across months learns small facts incrementally. In
March, the user mentions they prefer espresso. In August, they are cutting
caffeine and want tea in the morning. A naive memory system overwrites the
older preference. A less naive one stores both as parallel facts. Neither
resolves the conflict when the user later asks what to order. The policy
question for persistent memory is whether the system can keep both claims,
scope them, contest one with the other, and demote a durable claim when the
evidence shifts.

Most agent-memory benchmarks ask whether the system got the right answer. This
paper instead asks whether the memory policy had the right governable evidence
handle: the stored-claim identifier a policy uses to contest, scope, demote, or
repair memory. The gap shows up sharply in this paper's internal diagnostic
rows: B-cubed F1 is at or near 1.00 while exact policy-query canonical-id
agreement is 0.00. The system can group the right claims while failing to name
the stored claim a governance operation must address. Recall, retrieval, and
answer metrics remain important, but they can hide two different failures.
Upstream, aggregate extractor scores can authorize a candidate stream, the
shared extracted-fact list every policy receives, that collapses the policy
comparison. Downstream, memory-state metrics can expose relevant content while
failing to certify the governable evidence handle needed for contradiction,
scope, demotion, and repair. Recent memory benchmarks
broaden what systems are tested on. This paper narrows how policy-level credit
gets assigned.

The broader benchmark landscape makes this narrower question timely. STATE-Bench
evaluates memory in stateful enterprise workflows, EvoMemBench frames memory as
self-evolving state where no single memory form is uniformly sufficient, and
MemoryAgentBench includes FactConsolidation as an explicit revision and
consolidation target \cite{statebench2026,evomembench2026,memoryagentbench2025}.
This empirical evaluation-methods paper therefore argues for a stricter way to
assign credit, with committed artifacts for each empirical claim.
Memory-policy comparisons should hold the upstream candidate stream fixed and
use the same storage substrate, the shared memory store for compared policies.
They should separate oracle-mode policy claims (known scope, conflicts, and
sources) from noisy-pipeline claims (real extracted inputs). Noisy comparison
runs only after an extractor-quality gate, a check that the fact extractor is
adequate. Artifacts stay inspectable so failure cases can be audited. The running
case study is a lightweight consolidation queue policy
(CQ),\footnote{A non-archival design sketch is available at
\href{https://www.manavm.com/writing/the-consolidation-queue/}{The
Consolidation Queue: How Persistent AI Earns Durable Change}.} compared with
immediate-write Reflection, an ADD/UPDATE/NOOP eager-write baseline
(\memzero{}), implemented here as a narrow local comparator, and ablated
variants. We use CQ to test which memory-policy distinctions survive the
interfaces around them.

Concretely, CQ is a staged-write policy. Incoming candidate memories enter a
queue with scope, provenance, support, contradiction, and lifecycle metadata.
High-scoring candidates promote to durable memory. Lower-scoring but usable
candidates can remain pending and still answer scoped lookups. Contradictory
evidence contests earlier candidates and can demote active durable memories,
while narrower pending overrides can shadow a wider durable claim without
globally overwriting it. In the opening example, CQ keeps both the espresso
preference and the morning-tea preference as scoped claims: the narrower
morning-tea claim shadows the broader espresso claim at read time, and either
can be contested or demoted when stronger evidence arrives.

Internally, under a preregistered noisy policy comparison, where predictions and
gates were locked before running, CQ's
contradiction-edge-driven contestation and demotion mechanism remains visible
on forced contradiction and partially visible on preference drift. Other rows
stay recorded as nulls or PFLC failures, where the policy cannot query the
right stored claim. The same discipline also surfaces an
extractor false-unlock, where aggregate scores clear but policy comparison
should not be licensed. Externally, a controlled LongMemEval v1 probe tests
readout cardinality, meaning how many evidence sessions are returned. It
produces a stable negative result for base CQ on evidence completeness. A later
preregistered follow-up repairs that specific failure by
exposing all eligible same-slot pending evidence at answer time,
while a cardinality-capped Reflection control reproduces the original loss. The
paper preserves the failure and the repair as separate evidence.

\begin{table}[H]
\centering
\caption{Contribution map.}
\label{tab:contribution-map}
\small
\begin{tabular}{p{0.18\linewidth} p{0.72\linewidth}}
\toprule
Layer & Contribution\\
\midrule
Problem & Standard memory scores can pass while the governable evidence handle
is missing.\\
PFLC diagnostic & PFLC checks whether the governable evidence handle is exposed, with the
canonical-id diagnostic and formal counterexample showing the metric gap.\\
External replay & AMB LoCoMo replays third-party published contexts to separate
returned gold evidence ids, rank/context saturation, and answer residuals.\\
CQ case study & CQ is the worked policy used to test whether
same-input, same-store comparisons expose false unlocks, bounded nulls, and
interface failures.\\
Readout follow-up & Returning both pending evidence items, plus a capped
Reflection control, pins the LongMemEval v1 loss on answer-time readout
cardinality.\\
\bottomrule
\end{tabular}
\end{table}

The paper uses CQ as the case study for an attribution method. It is not a
benchmark leaderboard or a deployable memory-system release. The limits that
matter for interpretation are stated with the relevant results and in
Section~\ref{sec:limitations}.

\claim{Claim 1: Same-stream evaluation makes policy attribution inspectable}
In the main comparison, every policy receives the same ordered
\code{CandidateUpdate} stream, writes through the same scoped memory store, and
runs only after the extractor-quality gate records sufficient component
quality. That removes the main stream, substrate, and extractor-quality
confounds while preserving lifecycle traces, manifests, failure examples, and
replay checks.

\claim{Claim 2: Standard memory metrics can pass while the governable evidence handle is missing}
Cluster-partition, retrieval-at-k, and answer-accuracy metrics can all be
passed by systems whose policy cannot name the stored claim it must address.
PFLC asks whether the governable evidence handle is exposed, not just whether
retrieved text is semantically relevant.
The identity-PFLC gap is the memory-state counterpart of the extractor-gating
gap: the aggregate score looks good, but the governable evidence handle is
missing. The field survey shows when released benchmark artifacts do or do not
expose the per-question identifiers needed to replay such diagnostics.

\claim{Claim 3: A preserved negative result exposes interface failures}
The LongMemEval workstream adapts an external benchmark without rewriting away
the negative result. That can surface stable interface failures hidden by
answer-only, any-hit, or semantic-relevance scoring. The workstream builds a
redacted same-stream adapter, hidden-answer protocol, judge gate, PFLC scorer,
and contract-sensitivity audit before policy execution. The result is stable
and negative for base CQ on \code{all_hit_at_50}, which asks whether all
required evidence sessions are returned. Base CQ returns one required session,
but the missing second evidence session drives the loss.

\claim{Claim 4: A paired repair and control pin down readout cardinality}
The follow-up changes one read interface at a time. CQ is allowed to return all
eligible pending evidence in the same slot. Reflection keeps its winning write
path but is capped to one returned candidate at answer time. If the loss comes
from readout cardinality rather than storage asymmetry or input asymmetry, CQ
should close the gap and capped Reflection should lose in the same way. Both
predictions hold. This supports a local read-interface explanation for the
probe.

% Method-stage diagram for the four paper claims.
%
% The figure preserves the preregistered claim numbering while making the
% argument read as an attribution method rather than an evidence ledger:
% fair attribution -> lookup-contract diagnostic -> preserved loss plus paired
% repair/control.

\begin{figure}[H]
\centering
\resizebox{\linewidth}{!}{%
\begin{tikzpicture}[
  font=\small,
  stage/.style={
    rounded corners=3pt, draw, line width=0.8pt,
    minimum width=56mm, minimum height=38mm,
    align=center, inner sep=3mm,
  },
  s1/.style={stage, draw=blue!55!black, fill=blue!5},
  s2/.style={stage, draw=purple!55!black, fill=purple!5},
  s3/.style={stage, draw=orange!65!black, fill=orange!6},
  edge/.style={-{Stealth[length=2.6mm]}, line width=0.8pt},
  edgelbl/.style={
    font=\scriptsize, align=center,
    fill=white, draw=none, inner sep=1.5pt, text=black!70
  },
]

\node[s1] (stage1) {
  \begin{tabular}{@{}c@{}}
    {\footnotesize\bfseries\color{black!55} Stage 1}\\[-1pt]
    {\small\bfseries Fair attribution}\\[3pt]
    {\footnotesize\bfseries Claim 1}\\[-1pt]
    {\scriptsize Same stream}\\[-2pt]
    {\scriptsize Same substrate}\\[-2pt]
    {\scriptsize Extractor-gated unlock}\\[-2pt]
    {\scriptsize Inspectable traces}
  \end{tabular}
};

\node[s2, right=18mm of stage1] (stage2) {
  \begin{tabular}{@{}c@{}}
    {\footnotesize\bfseries\color{black!55} Stage 2}\\[-1pt]
    {\small\bfseries Lookup-contract diagnostic}\\[3pt]
    {\footnotesize\bfseries Claim 2}\\[-1pt]
    {\scriptsize PFLC asks for}\\[-2pt]
    {\scriptsize the policy-facing}\\[-2pt]
    {\scriptsize evidence handle}\\[-2pt]
    {\scriptsize not only relevant text}
  \end{tabular}
};

\node[s3, right=18mm of stage2] (stage3) {
  \begin{tabular}{@{}c@{}}
    {\footnotesize\bfseries\color{black!55} Stage 3}\\[-1pt]
    {\small\bfseries Preserved loss + paired repair}\\[3pt]
    {\footnotesize\bfseries Claim 3}\\[-1pt]
    {\scriptsize LongMemEval v1 loss}\\[-2pt]
    {\scriptsize preserved, not rescued}\\[2pt]
    {\footnotesize\bfseries Claim 4}\\[-1pt]
    {\scriptsize plural CQ repair}\\[-2pt]
    {\scriptsize capped Reflection control}
  \end{tabular}
};

\draw[edge] (stage1.east) -- node[edgelbl]
  {bounded nulls\\need handle tests} (stage2.west);
\draw[edge] (stage2.east) -- node[edgelbl]
  {evidence-id scoring\\on readout ids} (stage3.west);

\draw[edge, dashed, draw=black!55]
  (stage1.south) .. controls +(22mm,-14mm) and +(-22mm,-14mm) ..
  node[edgelbl, pos=0.5]
  {same-stream protocol\\carries to probe}
  (stage3.south);

\end{tikzpicture}%
}
\caption{Method-stage view of the four claims. Stage~1 makes policy attribution
inspectable by fixing the stream, substrate, and extractor gate. Stage~2 adds
PFLC to test whether the policy-facing evidence handle is exposed. Stage~3
preserves the LongMemEval v1 loss and then uses a paired repair/control to
isolate answer-time readout cardinality. The figure preserves the claim numbering but
frames the contribution as an attribution method.}
\label{fig:evidence-ledger}
\end{figure}


All empirical claims in this paper cite committed artifacts. Diagnostic-only,
descriptive-only, and artifact-blocked results are marked in the main text
rather than relegated to caveats.
